\section{Conclusion}
\label{sec:conclusion}

The cost of a zone-based TCTL tableau has two sources, and they need different
treatment. We measured which one dominates, found that on the most expensive
instances it is the discrete one and that the symbolic machinery contributes a
factor of exactly $1.00$, and removed a large part of the discrete factor with a
decomposition theorem imported in spirit from model counting. Conjunctions of
independent obligations go from $2^{k}$ to $O(k)$.

What we could not move is nesting. An $\AG$ inside an $\EF$ inside an $\AG$
shares atoms all the way down, decomposition does not apply, and eight further
attempts --- zone subsumption, $LU$ abstraction, forced stops, collapsing the
pursuit record, goal subsumption, encoding $\AG$ as context, and treating
single-polarity atoms as non-linking --- either failed outright or bought a few
per cent. Several failed instructively and are documented with the formula that
breaks them. Our present belief is that the remaining growth on those shapes is
the genuine $2^{O(n)}$ of the bound.

Two methodological points seem worth more than the speed-ups. First, the
profile was wrong three times running: the bottleneck was never where the
literature's emphasis on abstraction suggested, and each time it moved after
being fixed. Second, an evaluation suite that carries expected verdicts found a
soundness bug in the calculus that a regression suite of the usual kind had not,
and a regression check found a soundness bug in the decomposition that the
evaluation suite could not, because the suite had no data variables in it. The
two are complementary and we would not now build one without the other.
